% Clase 12 --- Qué se conserva y qué no al pasar del cable a un cuerpo
% cualquiera (§5.1 y §5.2).
% El docente dibujó "un conductor medio raro, con cualquier forma" justamente
% para mostrar que DeltaV = R I sigue valiendo pero R = rho L / A no. Los dos
% paneles tienen el mismo material: lo único que cambia es la forma y dónde se
% ponen los bornes.
\begin{center}
\begin{tikzpicture}[scale=0.95]

  % =============== (a) filiforme ===============
  \begin{scope}
    \fill[figblue!8] (-1.9,-0.55) rectangle (1.9,0.55);
    \draw[figgray, line width=1pt] (-1.9,-0.55) rectangle (1.9,0.55);

    % bornes
    \draw[figblue, line width=1.4pt] (-1.9,0) -- (-2.6,0);
    \draw[figblue, line width=1.4pt] (1.9,0) -- (2.6,0);
    \node[figcarga, fill=figblue, minimum size=5pt] at (-2.6,0) {};
    \node[figcarga, fill=figblue, minimum size=5pt] at (2.6,0) {};

    % campo y corriente, uniformes
    \foreach \y in {-0.28,0.28} { \draw[figvecres] (-1.4,\y) -- (1.4,\y); }
    \node[figlbl, figred, anchor=south] at (0,0.62) {$\vec E$ uniforme};

    % cotas
    \draw[figvecaux] (-1.9,-1.0) -- (-0.35,-1.0);
    \draw[figvecaux] (1.9,-1.0) -- (0.35,-1.0);
    \node[figlbl, anchor=north] at (0,-0.92) {$L$};
    \draw[figvecaux] (2.15,-0.55) -- (2.15,-0.15);
    \draw[figvecaux] (2.15,0.55) -- (2.15,0.15);
    \node[figlblaux, anchor=west] at (2.2,0.7) {sección $A$};

    \node[figlbl, figred, anchor=north, align=center] at (0,-1.95)
      {$R = \dfrac{\rho\,L}{A}$};
    \node[figlbl, anchor=north, align=center] at (0,-3.2)
      {(a) filiforme\\[-0.2em] \footnotesize toda la corriente recorre lo mismo};
  \end{scope}

  % =============== (b) forma cualquiera ===============
  \begin{scope}[xshift=8.4cm]
    \fill[figblue!8] (-2.0,-0.7) .. controls (-1.4,-1.3) and (0.4,-1.5) .. (1.9,-0.85)
      .. controls (2.5,-0.6) and (2.2,0.4) .. (1.5,0.75)
      .. controls (0.6,1.2) and (-1.2,1.05) .. (-1.9,0.55)
      .. controls (-2.35,0.2) and (-2.3,-0.4) .. (-2.0,-0.7);
    \draw[figgray, line width=1pt] (-2.0,-0.7) .. controls (-1.4,-1.3) and (0.4,-1.5) .. (1.9,-0.85)
      .. controls (2.5,-0.6) and (2.2,0.4) .. (1.5,0.75)
      .. controls (0.6,1.2) and (-1.2,1.05) .. (-1.9,0.55)
      .. controls (-2.35,0.2) and (-2.3,-0.4) .. (-2.0,-0.7);

    % par de bornes lejanos
    \node[figcarga, fill=figblue, minimum size=5.5pt] at (-1.95,0.15) {};
    \node[figcarga, fill=figblue, minimum size=5.5pt] at (2.25,-0.2) {};
    \draw[figblue, line width=0.9pt] (-1.95,0.15) .. controls (-0.6,0.5) and (0.9,-0.7) .. (2.25,-0.2);
    \node[figlbl, figblue, anchor=south, scale=0.9] at (0.15,0.55) {$R$ grande};

    % par de bornes cercanos
    \node[figcarga, fill=figred, minimum size=5.5pt] at (0.35,-1.32) {};
    \node[figcarga, fill=figred, minimum size=5.5pt] at (1.25,-1.24) {};
    \draw[figred, line width=0.9pt] (0.35,-1.32) .. controls (0.6,-1.0) and (1.0,-1.0) .. (1.25,-1.24);
    \node[figlbl, figred, anchor=north, scale=0.9] at (0.8,-1.45) {$R$ chica};

    \node[figlbl, figred, anchor=north, align=center] at (0,-1.95)
      {$R \neq \dfrac{\rho\,L}{A}$};
    \node[figlbl, anchor=north, align=center] at (0,-3.2)
      {(b) forma cualquiera\\[-0.2em] \footnotesize $R$ depende de dónde van los
       bornes};
  \end{scope}

\end{tikzpicture}\\[0.5em]
{\small\itshape Conviene separar las dos afirmaciones, porque una sobrevive al
cambio de forma y la otra no. Que $\Delta V$ sea proporcional a $I$ vale para
\textbf{cualquier} cuerpo de material óhmico, sea cual sea su forma y dónde se
conecte: siempre existe un número $R$. Lo que vale sólo en el caso filiforme es
la fórmula $R = \rho L/A$, porque ahí toda la corriente recorre el mismo largo
$L$ a través de la misma sección $A$ ---eso es lo que permite reemplazar $E$ por
$\Delta V/L$ y $J$ por $I/A$---. En el panel (b) no hay ningún $L$ ni ningún $A$
que poner: la corriente se reparte por caminos de distinto largo y sección.
Además $R$ deja de ser una propiedad del cuerpo solo y pasa a depender de dónde
se colocan los bornes, algo que se ve sin cuentas: bornes cercanos dejan a los
portadores un trayecto corto. La resistividad $\rho$, en cambio, no cambia
nunca: es del material.}
\end{center}
